\documentclass[12pt]{article}
\usepackage[T1]{fontenc}
\usepackage[utf8]{inputenc}
\usepackage{lmodern}
\usepackage{textcomp}
\usepackage{enumitem}
\usepackage{geometry}
\usepackage{graphicx}
\usepackage{amsmath, amssymb}
\geometry{margin=1in}
\begin{document}
\begin{center}
Economics - K-12 Level Assessment 15 \
Total Marks: 40 \
Answer all questions.
\end{center}

\section*{Multiple Choice (10 marks)}
\begin{enumerate}[label=\arabic*.]
\item The short-run Phillips curve indicates a
\begin{enumerate}[label=(\alph*)]
    \item direct relation between unemployment and inflation
    \item direct relation between price and quantity demanded
    \item inverse relation between price and quantity demanded
    \item inverse relation between unemployment and inflation
\end{enumerate}
\item Which of the following is true of a price floor?
\begin{enumerate}[label=(\alph*)]
    \item The price floor shifts the demand curve to the left.
    \item An effective floor creates a shortage of the good.
    \item The price floor shifts the supply curve of the good to the right.
    \item To be an effective floor, it must be set above the equilibrium price.
\end{enumerate}
\item The effect of the spending multiplier is lessened if
\begin{enumerate}[label=(\alph*)]
    \item the price level is constant with an increase in aggregate demand.
    \item the price level falls with an increase in aggregate supply.
    \item the price level is constant with an increase in long-run aggregate supply.
    \item the price level rises with an increase in aggregate demand.
\end{enumerate}
\item If the world price of copper exceeds the domestic (U.S.) price of copper we would expect
\begin{enumerate}[label=(\alph*)]
    \item the United States to be a net exporter of copper.
    \item the United States to impose a tariff on imported copper to protect domestic producers.
    \item the demand for U.S. copper to fall.
    \item a growing trade deficit in the United States in goods and services.
\end{enumerate}
\item The GDP Deflator differs from the CPI in that the GDP Deflator
\begin{enumerate}[label=(\alph*)]
    \item is thought to slightly overestimate the inflation rate
    \item uses base year quantities in its calculations
    \item incorporates both current year prices and base year prices
    \item incorporates current year quantities in its calculations
\end{enumerate}
\end{enumerate}

\section*{True / False (10 marks)}
\textit{Answer True or False}
\begin{enumerate}[label=\arabic*.]
\setcounter{enumi}{5}
\item True or False: Ceteris paribus means that anything can change at any time.
\item True or False: Automobiles provide the largest total utility and the smallest marginal utility due to their high demand and constant use.
\item True or False: Buying Treasury securities from commercial banks is a monetary tool used by the Federal Reserve.
\item True or False: An increase in aggregate demand will cause both unemployment and inflation rates to decrease at the same time.
\item True or False: National income measures the total income earned by the factors of production within a country.
\end{enumerate}

\section*{Long Form Response (20 marks)}
\textit{Answer each question with a clear and complete explanation.}
\begin{enumerate}[label=\arabic*.]
\setcounter{enumi}{10}
\item Discuss the role of firms in the circular flow model of a private closed economy, focusing on how they supply goods to households and the significance of this exchange in economic activity.
\item Discuss how the introduction of a new production technique that lowers costs can contribute to economic growth, providing examples and analyzing potential impacts on businesses and consumers.
\end{enumerate}

\vfill
\noindent\textit{End of Paper}
\end{document}